\documentclass[french]{book}
\usepackage{teach}
\usepackage{verbatim}
\usepackage{amsmath,amsfonts,amssymb}
\usepackage{pgfplots}
%\usepackage{geometry}
%\geometry{top=1cm, bottom=1cm, left=2cm, right=2cm}
\usepackage[utf8]{inputenc}
\usepackage[french]{babel}
\redefinecolor{thm}{title}{bgcolor}{alizarin}
\redefinecolor{prop}{title}{bgcolor}{meatbrown}
\redefinecolor{defn}{title}{bgcolor}{olivine}
\definecolor{alizarin}{rgb}{0.82, 0.1, 0.26}
\definecolor{meatbrown}{rgb}{0.9, 0.72, 0.23}
\definecolor{olivine}{rgb}{0.6, 0.73, 0.45}
\usepackage{pgf,tikz,tkz-tab}

\usepackage{mathrsfs}
\usetikzlibrary{arrows}

\begin{document}


\begin{center}
\huge{ Devoir surveillé IX}
\end{center}

\section*{Vecteurs et coefficient directeur}
\begin{enumerate}
\item Soit une droite (AB) tel que A(0;2) et B(1;5). Déterminer un vecteur directeur à la droite (AB) puis le coefficient directeur.
\item Soit une droite (OC) tel que son coefficient directeur soit $3,2$, décrire alors l'équation réduite de la droite.
\end{enumerate}

\section*{Appartenance d'un point à une droite}
Les points suivant appartiennent-ils à la droite $\Delta$ d'équation réduite $y=3x-5$? À la droite $\Delta_2$ ayant une équation cartésienne $4y+4x-2=0$? \textit{ (Pour récapituler vos résultats vous pouvez faire un tableau à 4 ligne et 2 colonne) }

\begin{enumerate}
\item A(-2;11)
\item B($\dfrac{3}{2}$;-1)
\item C(1,67;0)
\item D($\dfrac{7}{2}$;2)
\end{enumerate}

Pouvez vous donner le point d'intersection (si il existe) de la droite $\Delta$ et $\Delta_2$?

\section*{Déterminer une équation de droite}
Déterminer les équations de droite de: 
\begin{enumerate}
\item $\Delta_1$ dirigé par  $\overrightarrow{u} \left( \begin{array}{c} 2 \\ -2 \end{array} \right)$  Et passant par H(4;6)
\end{enumerate}


\section*{Associer une droite à l'équation d'une droite}
\begin{enumerate}
\item Donner une équation de droite est associée à $\Delta$
\item Tracer la droite d'équation réduite $\Delta_1: \; y = x-2$
\item Tracer la droite d'équation cartésienne  $\Delta_2: \; 2y+3x-1=0$
\item Tracer la droite d'équation cartésienne $\Delta_3: \; y-3=0$
\item Tracer la droite d'équation cartésienne $\Delta_4: \; x+5=0$
\end{enumerate}

\section*{Intersection de droite}
\begin{enumerate}
\item Soit deux droites $\Delta: \;  y=2x-1$ et $\Delta': \; y=3x+4$
\item Justifier à l'aide du coefficient directeur que les droites $\Delta$ et $\Delta'$ sont sécantes.
\item Déterminer alors le ou les point d'intersection de $\Delta$ et de $\Delta'$
\item (Bonus) Donner deux équations de droites différentes mais qui représente la même droite.
\end{enumerate}
\end{document}